\documentclass{article}
\usepackage{handover}
\usepackage{hyperref}


\begin{document}
\doctitle{Access and Credential Transfer}{Music Podcast Player}
\section{Purpose}

% Instructions for transfering control of the official 
% deployment and repository control. Mention that 
% instructions for deploying on AWS and Docker can be 
% found in the Deployment Instructions.

This document contains instructions for transferring 
control of the official GitHub repository and the official
AWS deployment of the Music Podcast Player between
accounts. Since this software is open-source, anyone may
deploy this software independently if they choose. For
instruction on on deployment, see ``Deployment 
Instructions.''

\section{GitHub Repository}
% Describes transfer of operational control of the 
% codebase.

This section discusses how to transfer operational 
control of the official Music Podcast Player GitHub
repository. This section will refer to the original
location of the repository as
\verb|https://github.com/{old-owner}/{old-name}| and the
new location of the repository as
\verb|https://github.com/{new-owner}/{new-name}|.

\subsection{Update OpenID Roles in AWS}
The Continuous Deployment actions in the repository rely
on OpenID to authenticate with AWS. In order to prevent
them from failing, the following AWS roles will need to
have their Trust Relationships updated.

\begin{enumerate}
    \item \verb|github-{environment}-deployment|
    \item \verb|github-{environment}-post-deploy|
\end{enumerate}

In each, the
\verb|token.actions.githubusercontent.com:sub| field of
each should be updated to include
\verb|repo:{new-owner}/{new-name}:environment:{environment}|.

The \verb|repo:{old-owner}/{old-name}:environment:{environment}|
entry can be kept until the second step is complete in
order to allow GitHub actions from the original repository
to continue working until the ownership is transferred. 
This entry should be removed once ownership has been
transferred.

\subsection{Transfer Ownership}

In the original repository, a user with admin permissions
or the owner must go to the repository settings, and under
General settings, select the button to transfer ownership
of the repository. Following this, the new owner must
accept the transfer.

Following the transfer, the new owner may review the roles
of collaborators under ``Collaborators and teams.''

\section{PodcastIndex API Credentials}

The Music Podcast Player depends on the Podcast Index API
for it's search feature. The Podcast Index requires
credentials to be supplied in order to complete an API 
request. These can be obtained from the
\href{https://api.podcastindex.org/}{Podcast Index API}
website.

In the event that these are rotated or changed for any
reason, they can be updated in an AWS deployment by 
storing the API key and API secret in the Systems Manager
Parameter Store, using the following parameters.

\begin{itemize}
    \item \verb|/rss-music-player/podcast-index-api/key|
    \item \verb|/rss-music-player/podcast-index-api/secret|
\end{itemize}

In a Docker deployment, they can be updated by storing the
new credentials in the following environment variables
via the \verb|secrets/backend.env| file.

\begin{itemize}
    \item \verb|RSS_PLAYER_PODCAST_INDEX_KEY|
    \item \verb|RSS_PLAYER_PODCAST_INDEX_SECRET|
\end{itemize}

\section{Official Deployment}
% Describes transfer of the official deployment in AWS.

% RDS snapshot and share.
% Transfer domain ownership.
% 
% 1. Deploy from CDK (includes DB, without wiring DNS) on account B 2. Write freeze and put site in maintenance mode 3. Snapshot database on account A 4. Restore database on account B 5. Test deployment on account B 6. Transfer R53 domain and point DNS to new deployment 7. Lift maintainance mode/write freeze.

This section describes the process for transferring the
official production deployment of the Music Podcast Player
to another AWS account. Note that this procedure is
designed to be as simple as possible while preserving user
data. If downtime is not acceptable, an alternative
procedure should be considered. Additionally, this
procedure is theoretical and has never been tested.

\subsection{Deploy on the New Account}

Begin by deploying to production on the new account but
without Route 53 Hosted Zone or A Record creation. The
deployment should have an empty database (aside from the
migrations run by the CDK) and a randomized domain though
Cloudfront. Email verification will also need to be
temporarily disabled, as this account does not have domain
access yet.

Currently, we support a flag to disable email
verification, but this works by bypassing email
verification for development rather than blocking those
features, which is what we'd recommend implementing for
this step.

Make sure to review ``Deployment Instructions'' and verify
the required Systems Manager Parameter Store secrets are
set up. (It's not a bad idea to rotate them either.)

\subsection{Write Lock and Database Snapshot}

The old account will then need to enable a write lock
mode, which will temporarily make the database read only
and write any cached changes. The key is the database
needs to be stable and not change until the write lock is
lifted.

We have not implemented a write lock or ``maintenance
mode,'' though we'd recommend one be implemented in the
future and in the event that the deployment needs to be
transferred.

Following the write lock taking effect, the old account
will need to manually make a snapshot of the database and
share access to it with the new account.

\subsection{Restore Snapshot to New Database}

In the new account, the snapshot of the database will be
restored to the new database. We'd recommend keeping the
snapshot for at least a week so the application can be
tested and time is given to discover data loss, should it
occur for any reason.

\subsection{Point DNS to New Deployment}

The DNS A record in the old account should be updated to
point to Cloudfront in the new account. The idea here is
that domain transfer can take a significant amount of time
and the site could be reopened to users (albeit without
email verification features enabled) during that time.
That said, this may not work because the site will not be
using the certificate for the domain until it has been
transferred, and, additionally, without the domain, the
CDK will point the frontend toward the Cloudfront URL for
the backend, likely causing CORS issues.

\subsection{Transfer Domain and Enable Email Verification}

Finally, transfer the domain from the old account to the
new one and update the deployment to use the official 
domain. Additionally, email verification can be enabled
at this point.

\subsection{Notes on this Procedure}

Since this has not yet been tested, we strongly recommend
testing this procedure with a ``production'' deployment in
two development environments before attempting to transfer
the official production deployment.

\end{document}
